\section{Research gaps}
\label{sec:research-gaps}
Three gaps motivate our work.

\emph{(i) Prediction circularity.} In a standard CBM the concepts are predicted
from the full image, so a concept can be inferred through global,
class-correlated shortcuts rather than from the evidence it
names~\cite{margeloiu2021concept}. The bottleneck then becomes circular: concept
and class predictions are entangled, undermining both the faithfulness of the
explanation and the reliability of concept interventions. Methods that ground
each concept on the specific body region it describes---preventing the concept
prediction from short-circuiting through the class---are still uncommon, and
closing this gap is the main motivation for our part-conditioned design.

\emph{(ii) Structure studied in isolation.} Structure in the concept space
(grounding, coarse-to-fine attributes) and in the label space (a taxonomy over
classes) have mostly been investigated separately; combining them within a single
supervised CBM remains under-explored.

\emph{(iii) Severity is not measured.} CBM evaluation almost always reports
accuracy alone, ignoring \emph{error severity}---the very quantity a class
taxonomy should improve. We address all three gaps by grounding concepts to
eliminate circularity, structuring both concept and label spaces, and evaluating
with an explicit severity metric. The backbone is held fixed across all
variants, so that observed differences are attributable to the bottleneck design
rather than to differences in feature extraction.
